\documentclass[a4paper]{article}
\usepackage[fontset=fandol]{ctex}
\usepackage[margin=1.65cm]{geometry}
\usepackage{amsmath,amssymb}
\usepackage{graphicx}
\usepackage{hyperref}
\usepackage{longtable,booktabs}
\usepackage{xcolor}
\usepackage{enumitem}
\usepackage{listings}
\lstset{basicstyle=\ttfamily\small,breaklines=true,columns=fullflexible,keepspaces=true}
\hypersetup{colorlinks=true,linkcolor=blue,urlcolor=blue}
\setlist[itemize]{leftmargin=1.25em,itemsep=0.18em}
\setlength{\parskip}{0.45em}
\setlength{\parindent}{2em}
\sloppy
\title{第 16 讲：RLVR：可验证奖励强化学习}
\author{CS336 Spring 2026 public videos and official course materials}
\date{}
\begin{document}
\maketitle
\tableofcontents
\newpage
\section{本章导读}
本章讨论 \textbf{RLVR：可验证奖励强化学习}。CS336 的第一层目标不是把现成大模型当作黑箱使用，而是把 language model 从输入文本、训练目标、模型结构、数据、系统和评估这些部件重新拆开。只有拆开之后，读者才能判断一个设计选择究竟改变了什么。
设想数学题或代码题可以自动判定正确性。RLVR（reinforcement learning from verifiable rewards）利用可验证奖励推动 reasoning，但也引入 verifier overfitting、length bias 和推理成本上升。
本章对应 CS336 Spring 2026 第 16 个公开视频，并结合课程页提供的可解析资料。视频和课程页给出事实边界；正文把这些材料改写为可自学的教材章节。
阅读本章时，先理解问题为什么存在，再看定义、公式和实现形态，随后通过例题和练习检查自己是否能独立复现这条 reasoning chain。
\begin{figure}[h]
\centering
\includegraphics[width=0.58\linewidth]{source_anchor.jpg}
\caption{第 16 讲的课程图像锚点。}
\end{figure}
\subsection{本章知识路线}
\begin{itemize}
\item 本章首先说明 RLVR：可验证奖励强化学习 为什么是 language modeling pipeline 中不可跳过的一环。
\item 其次定义本章反复使用的术语，并说明这些术语对应的计算对象或系统对象。
\item 随后用公式和伪代码表达主要机制，让概念能够落到可计算的变量上。
\item 最后给出例题、风险讨论和练习，便于读者检查自己是否真正掌握了机制。
\end{itemize}
\subsection{结构图}
下面的图用于帮助读者定位本章的核心结构。先看概念之间的依赖，再回到正文中的公式、伪代码和例题。
\begin{figure}[h]
\centering
\includegraphics[width=0.72\linewidth]{pdf_pages/page-02.jpg}
\caption{官方材料图页：RLVR：可验证奖励强化学习 的课程页/PDF/script 图像锚点。}
\end{figure}
\begin{figure}[h]
\centering

\includegraphics[width=0.72\linewidth]{pdf_pages/page-01.jpg}
\caption{官方材料图页：RLVR：可验证奖励强化学习 的课程页/PDF/script 图像锚点。}
\end{figure}
\section{本章主线}
本章可以沿着三条线索阅读：课程为什么把这个主题放在这里，它改变了哪些变量，以及怎样用小实验检查这些变量。
\subsection{本章要解决的核心问题}
本讲讲 RLVR：在数学、代码、可验证任务中，reward 可以由 verifier 给出，而不是依赖主观偏好标注。
这一点对应的学习动作是：找出对象，写出变量，说明变量如何被测量。若输入规模、硬件、数据分布或评估协议改变，结论也必须重新检查。
因此，RLVR：可验证奖励强化学习 不应被读成若干名词的集合，而应被读成一组可检验的机制。
\subsection{从机制到资源账本}
GRPO 等方法用组内相对 reward/advantage 训练 reasoning policy，减少 value function 复杂度，但也带来 length bias 和 baseline 合法性问题。
这一点对应的学习动作是：找出对象，写出变量，说明变量如何被测量。若输入规模、硬件、数据分布或评估协议改变，结论也必须重新检查。
因此，RLVR：可验证奖励强化学习 不应被读成若干名词的集合，而应被读成一组可检验的机制。
\subsection{从课堂讲解到可复现实验}
RLVR 的风险是过拟合 verifier、奖励稀疏、泛化不足和推理长度成本上升。评估必须同时看准确率、长度、失败类型和安全性。
这一点对应的学习动作是：找出对象，写出变量，说明变量如何被测量。若输入规模、硬件、数据分布或评估协议改变，结论也必须重新检查。
因此，RLVR：可验证奖励强化学习 不应被读成若干名词的集合，而应被读成一组可检验的机制。
\section{术语、符号与问题边界}
本章的术语横跨模型、数据和系统。读者需要先知道每个词指向的对象：有些词描述数学目标，有些词描述张量形状，有些词描述硬件瓶颈，还有一些词描述评估协议。术语边界不清，后面的公式和实现就会变成空洞的缩写。
\begin{longtable}{p{0.08\linewidth}p{0.42\linewidth}p{0.42\linewidth}}
\toprule
\textbf{\#} & \textbf{术语} & \textbf{阅读要求}\\
\midrule
1 & RLVR（reinforcement learning from verifiable rewards） & 说明它约束的是模型行为、优化目标还是评估协议，并指出一个奖励或偏好失真的例子。\\
2 & verifiable reward（可验证奖励） & 说明它约束的是模型行为、优化目标还是评估协议，并指出一个奖励或偏好失真的例子。\\
3 & GRPO & 给出定义、一个最小例子、一个关联公式或代码位置，以及一个典型失败情形。\\
4 & reasoning model（推理模型） & 给出定义、一个最小例子、一个关联公式或代码位置，以及一个典型失败情形。\\
5 & advantage（优势函数） & 给出定义、一个最小例子、一个关联公式或代码位置，以及一个典型失败情形。\\
\bottomrule
\end{longtable}
这些术语之间有依赖关系。例如 tokenization 改变 sequence length，sequence length 又改变 attention FLOPs、KV cache 和 evaluation token budget；parallelism 改变显存占用，也改变通信瓶颈和故障恢复方式。
\subsection{关键词}
下面的关键词在课程材料中反复出现。它们不代替定义，只提示读者本章哪些对象需要在后文持续跟踪。
\begin{itemize}
\item grpo, data, rlhf, rewards, reward, reasoning, kimi, qwen, expert, policy, practice, loss, length, level
\end{itemize}
\section{RLVR 的基本设置}
现在进入 RLVR 的基本设置。它把 RLVR：可验证奖励强化学习 中的一部分问题推进到可操作层面：哪些对象要被定义，哪些变量要被测量，哪些限制会改变结论。
\subsection{为什么要讨论 RLVR 的基本设置}
视频把本讲定位为第二节 post-training，聚焦 RLVR 和 reasoning 方向的及时进展。
RLVR 不依赖主观偏好 reward，而是利用数学题、代码测试、形式验证等可判定信号。
这使 reward 更客观，但任务覆盖范围更窄，并且容易过拟合到 verifier。
后训练部分的核心问题是行为控制。模型已经学到语言分布，但还需要通过示范、偏好或可验证奖励，把输出推向可用、可评估、可约束的行为。
\subsection{RLVR 的基本设置 的机制}
视频把本讲定位为第二节 post-training，聚焦 RLVR 和 reasoning 方向的及时进展。 这句话把讨论落在行为、奖励和偏好状态上。读者需要知道哪些量可直接观察，哪些量只能通过日志、profile 或评估协议间接推断。
RLVR 不依赖主观偏好 reward，而是利用数学题、代码测试、形式验证等可判定信号。 因而同一个方法在不同规模下可能改变不同的变量，尤其是reward 分布、偏好差异、KL 约束、回答长度和 verifier 通过率。比较方法时，必须同时给出问题规模和实验条件。
这使 reward 更客观，但任务覆盖范围更窄，并且容易过拟合到 verifier。 这使本节具有可检验性：一个结论若不能在公式、伪代码、profile trace、benchmark 或 ablation 中表现出来，就仍然只是直觉。
因此，RLVR 的基本设置 应当被理解为可以实现和测试的机制，而不是一个章节标题。CS336 反复强调 from scratch，正是因为只有能被复现的机制，才可能被稳定改进。
\subsection{RLVR 的基本设置 的目标函数}
\[R(x,y)=\mathbf{1}\{\operatorname{verify}(x,y)=\text{correct}\}\]
符号说明：\(x\) 是问题，\(y\) 是模型输出，verify 是外部可验证器；奖励常是稀疏二值信号。
这个关系式说明 RLVR 的基本设置 如何把行为偏好写入优化目标。关键不是符号复杂度，而是 reward、reference policy、采样分布和约束强度如何共同决定更新方向。
例如，validation loss 常被用作模型质量的 proxy，tokens/s 用作吞吐 proxy，Jaccard 或 MinHash 用作近重复 proxy，reward model score 用作偏好 proxy。proxy 不是目标本身；它只在评估协议清楚时才可引用。
\subsection{RLVR 的基本设置 的训练流程}
把 RLVR 的基本设置 写成训练流程时，要区分采样、打分、构造 loss 和更新参数。下面的伪代码保留这些边界。
\begin{lstlisting}[language=Python]
response = policy.generate(problem)
reward = verifier(problem, response)
update_policy_with_rl(reward)
\end{lstlisting}
这段伪代码的读法，是找出 RLVR 的基本设置 改变了哪些状态，以及这些状态如何对应到 reward 分布、偏好差异、KL 约束、回答长度和 verifier 通过率。如果某个状态没有被记录，后续实验就很难解释异常结果。
\subsection{边界条件与常见误解}
\begin{itemize}
\item 可验证不等于完整；很多真实任务没有可靠 verifier。
\item 稀疏 reward 会让探索和 credit assignment 变难。
\item 不要把小实验中的局部结论自动外推到 frontier-scale；scale、数据分布和硬件拓扑都会改变结论。
\item 不要把 benchmark 或 profile 的单个数字当作机制解释；必须同时记录实验设置、输入形状、版本和评估口径。
\end{itemize}
\subsection{与本章其他概念的关系}
RLVR 的基本设置 与本章其他部分的关系是互补的：术语表给出对象名称，公式给出最小数学结构，伪代码给出执行顺序，例题则把这些内容压缩成一个可检查的计算。
对这一类主题，最常见的误解是只列方法名。更有用的做法是比较每个方法改变的成本项、依赖的假设和可能的失败模式。
\subsection{怎样检验 RLVR 的基本设置}
检验 RLVR 的基本设置 时，先把抽象说法落到行为、奖励和偏好状态上。与 RLVR、verifiable、reward、math、code 相关的对象必须能被构造、打印、计数或评分；否则实验只能得到描述性观察，不能得到可复现结论。
一个合适的小实验应当保留本节机制，同时让 reward 分布、偏好差异、KL 约束、回答长度和 verifier 通过率 中至少一个量发生可解释变化。输入规模不必逼真；关键是读者能手算或打印关键中间量，并说明哪个变量触发了结果变化。
随后把公式说明转成可记录的数量关系。符号说明：\textbackslash{}(x\textbackslash{}) 是问题，\textbackslash{}(y\textbackslash{}) 是模型输出，verify 是外部可验证器；奖励常是稀疏二值信号。 如果数量只能间接观测，就必须说明使用了什么 proxy；如果使用 benchmark 或 reward，也要说明协议和随机性。
对照实验只应改变一个主要因素，例如 batch size、sequence length、vocabulary size、attention heads、filter threshold、learning rate、decoding 参数或 verifier。多个因素同时变化时，实验最多说明系统表现不同，不能说明是哪一个机制在起作用。
最后要主动寻找失败条件。教材中的成功路径容易记住，但工程中真正决定方法边界的是失败案例。本节至少要记住这些限制：可验证不等于完整；很多真实任务没有可靠 verifier；稀疏 reward 会让探索和 credit assignment 变难。复习时可以把这些限制改写成反例测试。
实验结论还要放回 RLVR：可验证奖励强化学习 的整体结构中。RLVR 的基本设置 会影响训练目标、数据选择、系统成本或评估解释中的至少一项；能说清楚这种影响，才说明读者已经从名词记忆进入机制理解。
\subsection{常见混淆：RLVR 的基本设置}
本节容易混淆的关键词包括 RLVR、verifiable、reward、math、code。它们在课程材料中可能连续出现，但层级并不相同：有的命名对象，有的描述操作，有的只是指标或约束。
其中，RLVR（reinforcement learning from verifiable rewards）使用可验证奖励训练推理能力；decode（解码）逐 token 生成，常受 memory bandwidth 和 KV cache 影响。
区分这些词时，可以先问它们是否改变 reward 分布、偏好差异、KL 约束、回答长度和 verifier 通过率。如果一个词不能对应到变量、状态、代码路径或评估协议，它在本节中就还没有形成可操作含义。
因此，复习 RLVR 的基本设置 时不应只抄关键词，而应写出至少一个最小例子：输入是什么，哪一步使用这个词，输出或指标怎样变化，失败条件在哪里出现。
\subsection{RLVR 的基本设置 的小结}
\begin{itemize}
\item 讨论 RLVR 的基本设置 时，应同时报告问题规模、实验设置和主要观测变量，尤其是 reward 分布、偏好差异、KL 约束、回答长度和 verifier 通过率。
\item 如果一个方法声称降低主要成本，要追问成本是否转移到了显存、通信、数据质量、训练稳定性或评估复杂度。
\item 公式和伪代码提供推理骨架；结论仍需要 profiler、ablation 或 held-out evaluation 支持。
\end{itemize}
\section{GRPO、优势估计与 length bias}
现在进入 GRPO、优势估计与 length bias。它把 RLVR：可验证奖励强化学习 中的一部分问题推进到可操作层面：哪些对象要被定义，哪些变量要被测量，哪些限制会改变结论。
\subsection{为什么要讨论 GRPO、优势估计与 length bias}
字幕关键词显示本讲反复讨论 GRPO、reward、advantage、length 和 DeepSeek。
group-relative 方法用同一 prompt 的多条采样构造相对优势，避免单独训练 value model 或降低复杂度。
reasoning RL 需要处理长输出：更长 chain-of-thought 可能提高正确率，也可能制造 length bias。
GRPO、优势估计与 length bias 的作用是把 RLVR：可验证奖励强化学习 中的一个抽象问题变成可计算对象：哪些变量进入公式，哪些状态进入代码，哪些条件会让结果失效。
\subsection{GRPO、优势估计与 length bias 的机制}
字幕关键词显示本讲反复讨论 GRPO、reward、advantage、length 和 DeepSeek。 这句话把讨论落在行为、奖励和偏好状态上。读者需要知道哪些量可直接观察，哪些量只能通过日志、profile 或评估协议间接推断。
group-relative 方法用同一 prompt 的多条采样构造相对优势，避免单独训练 value model 或降低复杂度。 因而同一个方法在不同规模下可能改变不同的变量，尤其是reward 分布、偏好差异、KL 约束、回答长度和 verifier 通过率。比较方法时，必须同时给出问题规模和实验条件。
reasoning RL 需要处理长输出：更长 chain-of-thought 可能提高正确率，也可能制造 length bias。 这使本节具有可检验性：一个结论若不能在公式、伪代码、profile trace、benchmark 或 ablation 中表现出来，就仍然只是直觉。
因此，GRPO、优势估计与 length bias 应当被理解为可以实现和测试的机制，而不是一个章节标题。CS336 反复强调 from scratch，正是因为只有能被复现的机制，才可能被稳定改进。
\subsection{GRPO、优势估计与 length bias 的数学表达}
\[A_i=\frac{R_i-\operatorname{mean}(R_{1:G})}{\operatorname{std}(R_{1:G})+\epsilon}\]
符号说明：同一问题采样 \(G\) 个回答，\(R_i\) 是第 \(i\) 个回答奖励，\(A_i\) 是组内标准化优势。
这个关系式说明 GRPO、优势估计与 length bias 如何把行为偏好写入优化目标。关键不是符号复杂度，而是 reward、reference policy、采样分布和约束强度如何共同决定更新方向。
例如，validation loss 常被用作模型质量的 proxy，tokens/s 用作吞吐 proxy，Jaccard 或 MinHash 用作近重复 proxy，reward model score 用作偏好 proxy。proxy 不是目标本身；它只在评估协议清楚时才可引用。
\subsection{从 GRPO、优势估计与 length bias 到程序}
把 GRPO、优势估计与 length bias 写成训练流程时，要区分采样、打分、构造 loss 和更新参数。下面的伪代码保留这些边界。
\begin{lstlisting}[language=Python]
samples = [policy.generate(x) for _ in range(G)]
rewards = [verify(x, y) for y in samples]
advantages = normalize_within_group(rewards)
policy_gradient_update(samples, advantages)
\end{lstlisting}
这段伪代码的读法，是找出 GRPO、优势估计与 length bias 改变了哪些状态，以及这些状态如何对应到 reward 分布、偏好差异、KL 约束、回答长度和 verifier 通过率。如果某个状态没有被记录，后续实验就很难解释异常结果。
\subsection{边界条件与常见误解}
\begin{itemize}
\item advantage normalization 会改变 reward scale，影响训练稳定性。
\item 如果 verifier 奖励格式技巧，模型可能学习投机输出。
\item 不要把小实验中的局部结论自动外推到 frontier-scale；scale、数据分布和硬件拓扑都会改变结论。
\item 不要把 benchmark 或 profile 的单个数字当作机制解释；必须同时记录实验设置、输入形状、版本和评估口径。
\end{itemize}
\subsection{与本章其他概念的关系}
GRPO、优势估计与 length bias 与本章其他部分的关系是互补的：术语表给出对象名称，公式给出最小数学结构，伪代码给出执行顺序，例题则把这些内容压缩成一个可检查的计算。
对这一类主题，最常见的误解是只列方法名。更有用的做法是比较每个方法改变的成本项、依赖的假设和可能的失败模式。
\subsection{怎样检验 GRPO、优势估计与 length bias}
检验 GRPO、优势估计与 length bias 时，先把抽象说法落到行为、奖励和偏好状态上。与 GRPO、advantage、DeepSeek、length、thinking 相关的对象必须能被构造、打印、计数或评分；否则实验只能得到描述性观察，不能得到可复现结论。
一个合适的小实验应当保留本节机制，同时让 reward 分布、偏好差异、KL 约束、回答长度和 verifier 通过率 中至少一个量发生可解释变化。输入规模不必逼真；关键是读者能手算或打印关键中间量，并说明哪个变量触发了结果变化。
随后把公式说明转成可记录的数量关系。符号说明：同一问题采样 \textbackslash{}(G\textbackslash{}) 个回答，\textbackslash{}(R\_i\textbackslash{}) 是第 \textbackslash{}(i\textbackslash{}) 个回答奖励，\textbackslash{}(A\_i\textbackslash{}) 是组内标准化优势。 如果数量只能间接观测，就必须说明使用了什么 proxy；如果使用 benchmark 或 reward，也要说明协议和随机性。
对照实验只应改变一个主要因素，例如 batch size、sequence length、vocabulary size、attention heads、filter threshold、learning rate、decoding 参数或 verifier。多个因素同时变化时，实验最多说明系统表现不同，不能说明是哪一个机制在起作用。
最后要主动寻找失败条件。教材中的成功路径容易记住，但工程中真正决定方法边界的是失败案例。本节至少要记住这些限制：advantage normalization 会改变 reward scale，影响训练稳定性；如果 verifier 奖励格式技巧，模型可能学习投机输出。复习时可以把这些限制改写成反例测试。
实验结论还要放回 RLVR：可验证奖励强化学习 的整体结构中。GRPO、优势估计与 length bias 会影响训练目标、数据选择、系统成本或评估解释中的至少一项；能说清楚这种影响，才说明读者已经从名词记忆进入机制理解。
\subsection{常见混淆：GRPO、优势估计与 length bias}
本节容易混淆的关键词包括 GRPO、advantage、DeepSeek、length、thinking。它们在课程材料中可能连续出现，但层级并不相同：有的命名对象，有的描述操作，有的只是指标或约束。
其中，GRPO 使用组内相对奖励/优势来训练 reasoning policy。
区分这些词时，可以先问它们是否改变 reward 分布、偏好差异、KL 约束、回答长度和 verifier 通过率。如果一个词不能对应到变量、状态、代码路径或评估协议，它在本节中就还没有形成可操作含义。
因此，复习 GRPO、优势估计与 length bias 时不应只抄关键词，而应写出至少一个最小例子：输入是什么，哪一步使用这个词，输出或指标怎样变化，失败条件在哪里出现。
\subsection{GRPO、优势估计与 length bias 的小结}
\begin{itemize}
\item 讨论 GRPO、优势估计与 length bias 时，应同时报告问题规模、实验设置和主要观测变量，尤其是 reward 分布、偏好差异、KL 约束、回答长度和 verifier 通过率。
\item 如果一个方法声称降低主要成本，要追问成本是否转移到了显存、通信、数据质量、训练稳定性或评估复杂度。
\item 公式和伪代码提供推理骨架；结论仍需要 profiler、ablation 或 held-out evaluation 支持。
\end{itemize}
\section{从 reasoning gains 到安全边界}
现在进入 从 reasoning gains 到安全边界。它把 RLVR：可验证奖励强化学习 中的一部分问题推进到可操作层面：哪些对象要被定义，哪些变量要被测量，哪些限制会改变结论。
\subsection{为什么要讨论 从 reasoning gains 到安全边界}
RLVR 的吸引力在于把模型推向更强 reasoning，但训练过程也可能增强规避、工具滥用或 reward hacking。
可验证任务上的提升不自动泛化到开放式推理、事实性或安全对话。
评估 RLVR 模型时要同时看 pass rate、过程长度、失败类型、泛化题集和安全行为。
从 reasoning gains 到安全边界 的作用是把 RLVR：可验证奖励强化学习 中的一个抽象问题变成可计算对象：哪些变量进入公式，哪些状态进入代码，哪些条件会让结果失效。
\subsection{从 reasoning gains 到安全边界 的机制}
RLVR 的吸引力在于把模型推向更强 reasoning，但训练过程也可能增强规避、工具滥用或 reward hacking。 这句话把讨论落在行为、奖励和偏好状态上。读者需要知道哪些量可直接观察，哪些量只能通过日志、profile 或评估协议间接推断。
可验证任务上的提升不自动泛化到开放式推理、事实性或安全对话。 因而同一个方法在不同规模下可能改变不同的变量，尤其是reward 分布、偏好差异、KL 约束、回答长度和 verifier 通过率。比较方法时，必须同时给出问题规模和实验条件。
评估 RLVR 模型时要同时看 pass rate、过程长度、失败类型、泛化题集和安全行为。 这使本节具有可检验性：一个结论若不能在公式、伪代码、profile trace、benchmark 或 ablation 中表现出来，就仍然只是直觉。
因此，从 reasoning gains 到安全边界 应当被理解为可以实现和测试的机制，而不是一个章节标题。CS336 反复强调 from scratch，正是因为只有能被复现的机制，才可能被稳定改进。
\subsection{从 reasoning gains 到安全边界 的数学表达}
\[\max_\theta\; \mathbb{E}_{y\sim\pi_\theta(\cdot\mid x)}[R(x,y)]-\beta\,\operatorname{KL}(\pi_\theta\|\pi_{\text{ref}})\]
符号说明：KL 项限制策略偏离参考模型，\(\beta\) 是约束强度。
这个关系式说明 从 reasoning gains 到安全边界 如何把行为偏好写入优化目标。关键不是符号复杂度，而是 reward、reference policy、采样分布和约束强度如何共同决定更新方向。
例如，validation loss 常被用作模型质量的 proxy，tokens/s 用作吞吐 proxy，Jaccard 或 MinHash 用作近重复 proxy，reward model score 用作偏好 proxy。proxy 不是目标本身；它只在评估协议清楚时才可引用。
\subsection{从 从 reasoning gains 到安全边界 到程序}
把 从 reasoning gains 到安全边界 写成训练流程时，要区分采样、打分、构造 loss 和更新参数。下面的伪代码保留这些边界。
\begin{lstlisting}[language=Python]
monitor = [accuracy, response_length, verifier_failures, safety_eval]
stop_if_reward_improves_but_generalization_drops(monitor)
\end{lstlisting}
这段伪代码的读法，是找出 从 reasoning gains 到安全边界 改变了哪些状态，以及这些状态如何对应到 reward 分布、偏好差异、KL 约束、回答长度和 verifier 通过率。如果某个状态没有被记录，后续实验就很难解释异常结果。
\subsection{边界条件与常见误解}
\begin{itemize}
\item 训练 reward 和真实目标之间的 gap 是核心风险。
\item reasoning 长度增加会显著提高推理成本。
\item 不要把小实验中的局部结论自动外推到 frontier-scale；scale、数据分布和硬件拓扑都会改变结论。
\item 不要把 benchmark 或 profile 的单个数字当作机制解释；必须同时记录实验设置、输入形状、版本和评估口径。
\end{itemize}
\subsection{与本章其他概念的关系}
从 reasoning gains 到安全边界 与本章其他部分的关系是互补的：术语表给出对象名称，公式给出最小数学结构，伪代码给出执行顺序，例题则把这些内容压缩成一个可检查的计算。
对这一类主题，最常见的误解是只列方法名。更有用的做法是比较每个方法改变的成本项、依赖的假设和可能的失败模式。
\subsection{怎样检验 从 reasoning gains 到安全边界}
检验 从 reasoning gains 到安全边界 时，先把抽象说法落到行为、奖励和偏好状态上。与 reasoning、KL、safety、generalization、reward hacking 相关的对象必须能被构造、打印、计数或评分；否则实验只能得到描述性观察，不能得到可复现结论。
一个合适的小实验应当保留本节机制，同时让 reward 分布、偏好差异、KL 约束、回答长度和 verifier 通过率 中至少一个量发生可解释变化。输入规模不必逼真；关键是读者能手算或打印关键中间量，并说明哪个变量触发了结果变化。
随后把公式说明转成可记录的数量关系。符号说明：KL 项限制策略偏离参考模型，\textbackslash{}(\textbackslash{}beta\textbackslash{}) 是约束强度。 如果数量只能间接观测，就必须说明使用了什么 proxy；如果使用 benchmark 或 reward，也要说明协议和随机性。
对照实验只应改变一个主要因素，例如 batch size、sequence length、vocabulary size、attention heads、filter threshold、learning rate、decoding 参数或 verifier。多个因素同时变化时，实验最多说明系统表现不同，不能说明是哪一个机制在起作用。
最后要主动寻找失败条件。教材中的成功路径容易记住，但工程中真正决定方法边界的是失败案例。本节至少要记住这些限制：训练 reward 和真实目标之间的 gap 是核心风险；reasoning 长度增加会显著提高推理成本。复习时可以把这些限制改写成反例测试。
实验结论还要放回 RLVR：可验证奖励强化学习 的整体结构中。从 reasoning gains 到安全边界 会影响训练目标、数据选择、系统成本或评估解释中的至少一项；能说清楚这种影响，才说明读者已经从名词记忆进入机制理解。
\subsection{常见混淆：从 reasoning gains 到安全边界}
本节容易混淆的关键词包括 reasoning、KL、safety、generalization、reward hacking。它们在课程材料中可能连续出现，但层级并不相同：有的命名对象，有的描述操作，有的只是指标或约束。
区分这些词时，可以先问它们是否改变 reward 分布、偏好差异、KL 约束、回答长度和 verifier 通过率。如果一个词不能对应到变量、状态、代码路径或评估协议，它在本节中就还没有形成可操作含义。
因此，复习 从 reasoning gains 到安全边界 时不应只抄关键词，而应写出至少一个最小例子：输入是什么，哪一步使用这个词，输出或指标怎样变化，失败条件在哪里出现。
\subsection{从 reasoning gains 到安全边界 的小结}
\begin{itemize}
\item 讨论 从 reasoning gains 到安全边界 时，应同时报告问题规模、实验设置和主要观测变量，尤其是 reward 分布、偏好差异、KL 约束、回答长度和 verifier 通过率。
\item 如果一个方法声称降低主要成本，要追问成本是否转移到了显存、通信、数据质量、训练稳定性或评估复杂度。
\item 公式和伪代码提供推理骨架；结论仍需要 profiler、ablation 或 held-out evaluation 支持。
\end{itemize}
\section{综合讨论：从局部机制到完整系统}
本章的主题是 RLVR：可验证奖励强化学习，但它不应被读成几个相邻知识点的拼接。RLVR 的基本设置、GRPO、优势估计与 length bias、从 reasoning gains 到安全边界 共同回答同一个问题：语言模型系统中哪些对象可以被定义，哪些成本可以被估算，哪些结论必须通过实验或评估来确认。
这些对象包括 RLVR（reinforcement learning from verifiable rewards）、verifiable reward（可验证奖励）、GRPO、reasoning model（推理模型）、advantage（优势函数）。读者可以把它们看成一组接口：有的接口连接文本和 token，有的连接模型结构和张量计算，有的连接数据样本和训练目标，有的连接离线评估和在线服务。接口一旦改变，后面的成本、错误类型和优化空间也会随之改变。
\subsection{RLVR 的基本设置 在系统中的位置}
RLVR 的基本设置 首先提供一个局部解释：视频把本讲定位为第二节 post-training，聚焦 RLVR 和 reasoning 方向的及时进展。 这类解释的价值在于把直觉压缩成可操作对象。一旦对象明确，读者就可以追问它的输入、输出、中间状态和失败条件，而不是只记住方法名。
其次，RLVR 的基本设置 会改变至少一个资源或评估变量。它可能改变 sequence length、activation size、memory bandwidth、communication volume、data mixture、reward distribution 或 benchmark protocol。这些变量在小例子里可以手算，在真实训练和服务系统里则需要 profiling、logging 和 held-out evaluation。
最后，RLVR 的基本设置 不能脱离边界条件理解。一个关键 caveat 是：可验证不等于完整；很多真实任务没有可靠 verifier。。把 caveat 写出来不是为了削弱结论，而是为了说明结论在哪些输入、硬件、数据和评估条件下才成立。
\subsection{GRPO、优势估计与 length bias 在系统中的位置}
GRPO、优势估计与 length bias 首先提供一个局部解释：字幕关键词显示本讲反复讨论 GRPO、reward、advantage、length 和 DeepSeek。 这类解释的价值在于把直觉压缩成可操作对象。一旦对象明确，读者就可以追问它的输入、输出、中间状态和失败条件，而不是只记住方法名。
其次，GRPO、优势估计与 length bias 会改变至少一个资源或评估变量。它可能改变 sequence length、activation size、memory bandwidth、communication volume、data mixture、reward distribution 或 benchmark protocol。这些变量在小例子里可以手算，在真实训练和服务系统里则需要 profiling、logging 和 held-out evaluation。
最后，GRPO、优势估计与 length bias 不能脱离边界条件理解。一个关键 caveat 是：advantage normalization 会改变 reward scale，影响训练稳定性。。把 caveat 写出来不是为了削弱结论，而是为了说明结论在哪些输入、硬件、数据和评估条件下才成立。
\subsection{从 reasoning gains 到安全边界 在系统中的位置}
从 reasoning gains 到安全边界 首先提供一个局部解释：RLVR 的吸引力在于把模型推向更强 reasoning，但训练过程也可能增强规避、工具滥用或 reward hacking。 这类解释的价值在于把直觉压缩成可操作对象。一旦对象明确，读者就可以追问它的输入、输出、中间状态和失败条件，而不是只记住方法名。
其次，从 reasoning gains 到安全边界 会改变至少一个资源或评估变量。它可能改变 sequence length、activation size、memory bandwidth、communication volume、data mixture、reward distribution 或 benchmark protocol。这些变量在小例子里可以手算，在真实训练和服务系统里则需要 profiling、logging 和 held-out evaluation。
最后，从 reasoning gains 到安全边界 不能脱离边界条件理解。一个关键 caveat 是：训练 reward 和真实目标之间的 gap 是核心风险。。把 caveat 写出来不是为了削弱结论，而是为了说明结论在哪些输入、硬件、数据和评估条件下才成立。
\subsection{把本章知识用于读论文和复现实验}
读相关论文时，可以把本章变成一个检查表。第一，论文是否清楚定义了输入规模、模型规模、数据来源和评估协议；第二，论文声称的改进落在哪个变量上；第三，作者是否报告了会让结论失效的设置；第四，实验是否能分离模型结构、数据质量、优化设置和系统实现的影响。
复现实验时，也应避免直接追求大规模。先用最小程序复现一个公式、一个伪代码分支或一个 profile 现象，再扩大输入规模。这样做的原因很简单：如果小规模程序无法解释中间量，大规模训练只会把错误隐藏在日志和随机波动里。
把这些检查合在一起，本章给出的不是一个固定答案，而是一种阅读 CS336 的方法：每个模型机制都要落到变量，每个变量都要落到测量，每个测量都要说明协议，每个协议都要承认边界。
因此，本章中的任何一句结论都可以被改写成一个可引用的完整句式：在给定数据、模型、硬件和评估协议下，某个机制改变了某个变量，并通过某个观测指标表现出来。这样的句子比单独背诵术语更长，但它包含了自学、复现和引用时真正需要的信息。
把本章用于自学时，可以按三遍阅读。第一遍只追踪对象和定义，确认每个术语到底指向文本、token、张量、样本、请求还是评估记录。第二遍追踪公式和伪代码，确认每一步改变了哪些状态。第三遍追踪实验和 caveat，确认哪些结论只在特定规模、数据或硬件条件下成立。
把本章用于复习时，则应反过来操作：先从一个失败案例或异常指标出发，再回到相应机制。loss 曲线异常可能指向数据、优化或模型容量；吞吐异常可能指向 kernel、通信或调度；benchmark 异常可能指向 contamination、prompt protocol 或采样设置。能够从异常反推机制，是掌握本章的实用标准。
\section{例题：把概念变成可计算检查}
\subsection{GRPO 的组内优势}
同一数学题采样 8 个回答，用 verifier 给出 0/1 奖励。
\begin{itemize}
\item 计算组内平均和标准差，把每个回答的 reward 标准化成 advantage。
\item 正确回答 advantage 为正，错误回答 advantage 为负。
\item 如果所有回答都错或都对，学习信号会变弱。
\end{itemize}
\textbf{结论。} RLVR 的训练信号质量取决于采样多样性、verifier 可靠性和奖励分布。
这个例题的目的不是替代完整工程实现，而是给出一种最小推理形式：把问题转成变量，写出近似公式或伪代码，再说明近似何时失效。
\section{实验设计：从小例子到可复现结论}
学习 RLVR：可验证奖励强化学习 不能停在概念层。一个教材化结论至少需要四件事：可构造的输入、可观测的中间量、可比较的指标，以及清楚的失败条件。下面的实验设计不是要求训练大模型，而是要求读者用小规模程序复现本章机制。
\subsection{最小输入}
先构造只触发本章核心机制的小输入。例如 tokenizer 章节可以用几个包含 rare words 和 Unicode 字符的字符串；GPU 章节可以用一个 elementwise kernel 和一个 GEMM；data 章节可以用十几篇近重复文档。小输入的价值在于它让错误可见。
\subsection{资源账本}
对同一输入写下 FLOPs、bytes moved、显存峰值、通信量或样本数量的估算。估算不需要一开始就精确，但必须说明单位和近似假设。随后用 profiler、benchmark、validation loss 或人工检查去验证估算是否同量级。
\subsection{单变量消融}
一次只改变一个因素，例如 vocabulary size、head 数、batch size、filter threshold、reward scale、decoding temperature 或 GPU 数量。若多个变量同时变化，实验只能说明系统变了，不能说明机制是什么。
\subsection{失败条件}
最后要刻意触发失败：极端 context length、错误 dtype、过小 batch、污染 benchmark、低质量数据或不可靠 verifier。失败条件常常比成功曲线更能说明一个方法的边界。
\subsection{记录与报告}
实验记录应能让另一个读者复现同一结论。最少要写清楚输入构造、代码版本、随机种子、硬件、batch shape、数据版本、评价指标和异常处理方式。报告结论时不要只写“更快”或“更好”，而要说明快在哪里、好在哪个指标上，以及是否牺牲了内存、稳定性、数据质量或可解释性。
常见报告错误是把局部现象写成普遍规律。例如一次小模型实验里的 loss 改善，不等于更大模型或不同数据分布上也会改善；一次 kernel benchmark 的吞吐提升，不等于端到端 serving latency 一定降低。教材中的实验报告应始终把结论限定在可复现条件内。
\subsection{从小实验到大结论}
小实验的作用不是替代课程中的完整训练或系统实现，而是建立一条可检查的推理链。先在小输入上确认变量方向，再在中等规模上确认数量级，最后才讨论是否能外推到更大模型、更多 token 或更复杂 workload。缺少中间层级时，大结论通常只是在复述直觉。
因此，实验报告最好同时包含正例和反例：正例说明机制在受控条件下怎样工作，反例说明条件稍变时哪里失效。对 CS336 这样的课程，反例尤其重要，因为 language model 的错误常来自多个层面叠加：数据分布、优化稳定性、硬件瓶颈、评估协议和服务负载都可能同时改变观察结果。
还要记录资料版本与复现边界。公开视频、课程页、配套代码和 PDF 可能在不同时间更新；复现实验应写明使用的材料版本，并说明哪些结论来自课程事实，哪些是为了帮助理解而加入的解释性推导。这样才能把学习笔记变成可检查的技术记录。
若复现实验与课程结论不一致，首先检查输入、版本和评估协议，而不是立即修改模型假设。
\section{课程资料与回看路径}
教材正文已经把主要概念、公式和实现逻辑组织成连续叙述；本节只提供回到课程材料的阅读路径。读者复习时可以先完成前文练习，再按下面的时间段或资料组核对细节。
\subsection{视频回看路径}
\begin{longtable}{p{0.20\linewidth}p{0.46\linewidth}p{0.26\linewidth}}
\toprule
\textbf{时间} & \textbf{关键词} & \textbf{复习任务}\\
\midrule
00:00:05-00:06:23 & going, very, we're, that's, reinforcement, openai, we've & 回看定义、例子、推导或 caveat 在该段中的位置。\\
00:18:09-00:24:19 & grpo, standard, it's, well, deviation, baseline, going & 回看定义、例子、推导或 caveat 在该段中的位置。\\
00:37:05-00:43:31 & very, think, quite, it's, reward, data, talk & 回看定义、例子、推导或 caveat 在该段中的位置。\\
00:50:04-00:56:32 & very, training, then, your, inference, think, qwen & 回看定义、例子、推导或 caveat 在该段中的位置。\\
01:09:42-01:15:30 & long, then, it's, yeah., training, very, data & 回看定义、例子、推导或 caveat 在该段中的位置。\\
\bottomrule
\end{longtable}
这些时间段用于定位视频中的讲解顺序。严肃引用时，应回到视频片段确认讲者表述、板书或幻灯片内容。
回看视频时，不必逐字抄录字幕。更有效的方法是把每个时间段判定为定义、例子、推导、代码解释或 caveat，再把它放回本章对应小节。
\subsection{课程资料阅读路径}
\begin{longtable}{p{0.14\linewidth}p{0.48\linewidth}p{0.30\linewidth}}
\toprule
\textbf{资料组} & \textbf{标题/页块} & \textbf{关键词}\\
\midrule
official\_01 & Lecture 16 P OST- TRAININ G 2 R E INF O R C E M E NT L E A R NIN G F R O M V ER I F ... & practice, policy, reward, grpo, level, rewards, implementation\\
official\_03 & Interesting phenomena (?) Longer CoTs during training ‘aha’ moment? & reasoning, rlhf, grpo, language, non-verifiable, data, step\\
official\_04 & How well does R1 work? It’s pretty good (but you knew that) & kimi, distillation, long, policy, curation, reward, length\\
official\_05 & Kimi’s RL setup in detail & data, results, qwen, kimi, scaling, interesting, comes\\
official\_07 & Agent RL & agent, recap, overoptimization, problem, narrow, domains, solution\\
\bottomrule
\end{longtable}
课程资料通常给出更稳定的标题、公式、图或代码结构；视频则补充动机和口头解释。两者合起来构成本章的事实边界。
阅读资料时，可以把每个资料组拆成三个对象：一个定义，一个公式或伪代码动作，一个边界条件。这样比逐字翻译页面或脚本更容易形成可用知识。
\section{实践说明、历史位置与风险}
RLVR：可验证奖励强化学习 在 CS336 的课程结构中连接基础机制和规模化问题。基础机制解释方法为什么成立；规模化问题检验它在更大模型、更长上下文、更复杂数据或真实 serving workload 下是否仍成立。
历史位置不等于模型年表。更重要的是识别问题如何迁移：早期方法解决小规模建模问题，现代方法常常解决同一问题在大规模数据、GPU 集群、长上下文或人类偏好约束下的新形态。
\begin{itemize}
\item 资料版本限制（version risk）：课程页、公开视频和配套资料可能在学期中更新；引用时应注明访问日期或资料版本。
\item 测量口径限制（measurement risk）：FLOPs、tokens/s、benchmark score、reward 等数字只有在协议完整时可比较。
\item 规模外推限制（scaling risk）：小模型或小 batch 的机制不必然外推到 frontier-scale。
\item 数据分布限制（data risk）：数据来源、过滤和去重策略会改变模型行为，也会改变 evaluation contamination 风险。
\item 对齐风险（alignment risk）：post-training 或 multimodal tuning 的行为变化可能提高 helpfulness，同时降低校准、透明度或安全边界。
\end{itemize}
\section{复习要点}
下面的问题把本章内容压缩为可反复检查的条目。每个条目都应能回到前文的解释、公式、伪代码或例题。
\begin{longtable}{p{0.16\linewidth}p{0.38\linewidth}p{0.38\linewidth}}
\toprule
\textbf{条目} & \textbf{对象} & \textbf{检查动作}\\
\midrule
条目 1 & RLVR（reinforcement learning from verifiable rewards） & 定义它；写出一个最小例子；说明一个 failure mode；指出它影响哪个资源或评估账本。\\
条目 2 & verifiable reward（可验证奖励） & 定义它；写出一个最小例子；说明一个 failure mode；指出它影响哪个资源或评估账本。\\
条目 3 & GRPO & 定义它；写出一个最小例子；说明一个 failure mode；指出它影响哪个资源或评估账本。\\
条目 4 & reasoning model（推理模型） & 定义它；写出一个最小例子；说明一个 failure mode；指出它影响哪个资源或评估账本。\\
条目 5 & advantage（优势函数） & 定义它；写出一个最小例子；说明一个 failure mode；指出它影响哪个资源或评估账本。\\
\bottomrule
\end{longtable}
\subsection{综合自测}
\begin{itemize}
\item 我能否不用课程原句，独立解释本章中心问题？
\item 我能否把本章至少一个公式改写成可运行的估算函数？
\item 我能否指出一个看似改进但实际转移成本的方法？
\item 我能否回到视频和课程资料，找到支持本章关键论断的位置？
\item 我能否设计一个最小 ablation，验证本章一个 caveat？
\end{itemize}
\section{总结与延伸}
\subsection{本章小结}
\begin{itemize}
\item RLVR：可验证奖励强化学习 的核心不是记忆术语，而是理解对象、变量和机制之间的关系。
\item 视频给出讲解顺序，课程资料提供可核对的公式、图、代码或页面结构。
\item 复习时应把每个术语连接到至少一个变量、一个实现动作和一个失败模式。
\end{itemize}
\subsection{延伸解释}
本章中的延伸解释用于连接课程材料中的概念，例如说明公式里的 proxy、解释系统 tradeoff，或把多个材料片段串成一个完整推理链。若需要严肃引用，应回到课程视频和配套资料核对原始表述。
\subsection{练习}
\begin{itemize}
\item 定义题：用中英双语解释本章任意五个重要术语，并说明它们属于 compute、memory、data、optimization、evaluation 还是 deployment 账本。
\item 推导题：选择本章一个公式，逐个解释符号、单位和近似假设，并说明哪个变量最难在真实系统中测量。
\item 实现题：把本章伪代码改写成一个最小 Python sanity check，不要求训练大模型，但要输出可检查的中间量。
\item 资料题：从“课程资料与回看路径”中选择一个视频时间段，回到原始视频核对讲者例子，并记录是否存在字幕误差。
\item 评估题：为本章方法设计一个 ablation 或 benchmark，说明控制变量、数据版本、指标和 failure criteria。
\item 边界题：说明本章一个结论在更长 context、更大 batch、更多 GPU 或更复杂数据分布下可能如何失效。
\item 综合题：把本章内容和前一章或后一章连接起来，写出一个端到端训练/推理 pipeline 中的依赖关系。
\item 批判题：指出一个容易被营销材料夸大的说法，并用本章的资源账本或评估协议反驳它。
\end{itemize}
\end{document}
